\section{Fuente y pregunta central: detrás de un buen modelo no hay un solo comando de entrenamiento}

La versión anterior no incluía la canonical video URL, solo repetía la YouTube thumbnail como única figura de la clase y reducía el deck oficial de 71 páginas a una descripción breve. Esta clase se basa ahora en las Google Slides enlazadas directamente desde el sitio web personal de la ponente, en la grabación oficial de Stanford Online \texttt{jm2hyJLFfN8}, en los subtítulos oficiales \texttt{en-US} hechos por personas y en las primary sources de BigCode/StarCoder. Las 71 páginas de slides ya están guardadas en disco; 58 de ellas tienen contenido didáctico propio y las demás son solo separadores, estados intermedios de animaciones, marcas de finalización, páginas repetidas o agradecimientos.

\subsection{La clase no se reduce a la frase «los datos son lo más importante», sino que es una cadena de preguntas}

El título presenta StarCoder como use case, lo que significa que el objetivo de la clase no es repetir la configuración de un modelo, sino usar un modelo de código para responder preguntas generales de preentrenamiento: cuántos datos se necesitan, dónde encontrarlos, cómo filtrarlos, cómo gobernarlos, cómo evaluar y cuándo confiar en un benchmark. Cada capítulo posterior se desarrolla siguiendo esta estructura question--answer.

\lecturefigure{slide-01.jpg}{Behind the Scenes of LLM Pretraining: StarCoder use case}{Deck oficial de Loubna Ben Allal, página 1}

\lecturefigure{slide-02.jpg}{Ruta de la clase: responder «cómo entrenar un buen LLM» con una serie de preguntas}{Deck oficial de Loubna Ben Allal, página 2}

\begin{importantbox}{Las seis preguntas que esta clase debe responder}
\begin{enumerate}
  \item Con un compute dado, ¿cómo se reparten los parámetros del modelo y los tokens de entrenamiento?
  \item ¿Cómo se convierten las raw web/code/synthetic sources en un corpus entrenable?
  \item ¿Cómo se valida la intuición detrás de un filter mediante ablation con modelos pequeños?
  \item ¿Cómo entran la license, el opt-out, la PII y la decontamination al pipeline?
  \item ¿Cómo se da formato de sequence de entrenamiento a un repository/file/commit/notebook?
  \item ¿Cómo evita un leaderboard la contamination, la dependencia de un solo benchmark y la fuga temporal?
\end{enumerate}
\end{importantbox}

\teachervoice{00:00:50--00:02:10, Ben Allal dice que toda la exposición tiene una sola pregunta, «muy loaded», y que las slides irán planteando subpreguntas para luego responderlas. StarCoder es el caso que lleva los principios abstractos a un proyecto real.}

\subsection{Avances de los Open-model: qué capacidades prácticas trae abrir los pesos}

El memo ``No Moat'', que circuló en 2023, representa un juicio: que los modelos cerrados vayan adelante no significa que los abiertos nunca puedan alcanzarlos. La clase no lo cita para declarar terminada la competencia, sino para explicar por qué publicar los weights genera un nuevo ecosistema de deployment/fine-tuning.

\lecturefigure{slide-03.jpg}{El juicio «No moat»: la comunidad abierta completa poco a poco las piezas del entrenamiento}{Deck oficial de Loubna Ben Allal, página 3}

Una vez abiertos los weights, la quantization puede reducir el memory footprint, un runtime local permite desplegar en consumer hardware y el parameter-efficient fine-tuning permite personalizar con los datos de una organización. Este valor no equivale a que el modelo alcance la capacidad de los sistemas cerrados en todas las tareas, sino a que \textbf{el espacio de opciones verificables, modificables y desplegables de forma privada es mayor}.

\lecturefigure{slide-04.jpg}{Ejemplo de la clase de 2024: un modelo con evaluaciones cercanas a GPT-4 puede ejecutarse en un equipo de escritorio de consumo de gama alta}{Deck oficial de Loubna Ben Allal, página 4}

\begin{warningbox}{Un «rendimiento cercano» debe ir ligado a una versión, una métrica y una fecha}
La slide es una instantánea de mayo de 2024. Arena, MMLU, coding, seguridad, contexto largo, tool use, latency y cost dan clasificaciones distintas; esta clase no convierte la cercanía en una tabla concreta en una parity general.
\end{warningbox}

\teachervoice{00:02:10--00:03:20, la ponente enfatiza que el beneficio directo de los open weights es la quantization, el consumer deployment y el fine-tuning comunitario, no una frase abstracta como «lo abierto es mejor».}

\subsection{Las open releases aumentan, pero la transparencia del entrenamiento sigue siendo insuficiente}

Modelos públicos como Gemma y Mistral amplían el open ecosystem; la gráfica de Arena también muestra que la open/closed gap se redujo entre 2023 y 2024. Al leer estas dos páginas, primero hay que ver qué pregunta responde cada una: la primera trata del crecimiento de la oferta; la segunda, de la tendencia bajo un protocolo específico de evaluación de preferencias humanas.

\lecturefigure{slide-05.jpg}{Aumentan los lanzamientos de modelos abiertos: Gemma, Mistral y otros entran a la competencia}{Deck oficial de Loubna Ben Allal, página 5}

\lecturefigure{slide-06.jpg}{Perspectiva de Arena en May 2024: se reduce la brecha entre modelos open/closed}{Deck oficial de Loubna Ben Allal, página 6}

\begin{knowledgebox}{Qué puede demostrar la evidencia de Arena}
Arena forma un Elo-like ranking a partir de la pairwise preference de los usuarios, lo que sirve para medir la preferencia por respuestas interactivas; por sí solo no puede medir la legalidad de los datos de entrenamiento, la confiabilidad factual, la rare-language coverage, la code execution correctness ni el costo de despliegue. La tendencia sirve como referencia, pero los límites de capacidad siguen requiriendo una evaluación multidimensional.
\end{knowledgebox}

\subsection{Por qué se publican los pesos pero no los datos ni la recipe}

La ponente da dos razones prácticas: primero, la disclosure puede provocar license/copyright scrutiny; segundo, los equipos quieren conservar su competitive edge. Así, los weights del modelo pueden descargarse, pero no es posible saber de dónde vienen las muestras, cómo se eligieron los filters, cómo se asignó el número de tokens ni si hubo fuga en la evaluation.

\lecturefigure{slide-07.jpg}{Limitaciones de los modelos públicos: los datos y los detalles de entrenamiento suelen faltar}{Deck oficial de Loubna Ben Allal, página 7}

\begin{warningbox}{La falta de reproducibilidad no es «un poco menos de documentación»}
Si faltan la dataset version, el dedup threshold, los sampling weights, el tokenizer, los training tokens, el optimizer schedule, la benchmark decontamination y las seeds, los investigadores no pueden determinar si el rendimiento proviene de la architecture, de los data o de una fuga en la evaluación, ni pueden reutilizar de forma confiable las decisiones de gobernanza.
\end{warningbox}

\teachervoice{00:04:10--00:05:05, la ponente expone abiertamente las razones legales y competitivas para no divulgar. No por ello renuncia a resumir regularidades públicas, sino que enfatiza la necesidad de reunir evidence a partir de varias releases públicas.}

\subsection{Model, GPUs, Data: de los tres vértices del triángulo, esta clase elige data}

Transformer, Mamba y MoE pueden ser opciones de model, y las GPU fijan el límite del budget; pero con architecture y training technique similares, la calidad y la mixture de los datos suelen ser la principal diferencia dentro de un presupuesto fijo. Aquí, el «data backbone» es el focus de la clase; no se afirma que la architecture y el hardware sean irrelevantes.

\lecturefigure{slide-11.jpg}{Los tres tipos de recursos para entrenar un buen LLM: Model, GPUs y Data}{Deck oficial de Loubna Ben Allal, página 11}

\begin{knowledgebox}{Escribir los tres como una optimización con restricciones}
El objetivo de entrenamiento puede escribirse de forma aproximada como
\[
\min_{\theta,\mathcal{D},r}\ L(\theta;\mathcal{D},r)
\quad\text{s.t.}\quad
C_{\text{train}}\le B,\ C_{\text{serve}}\le S,
\]
donde $\mathcal{D}$ son los datos y la mixture, $r$ es la recipe de entrenamiento y $B,S$ son, respectivamente, los presupuestos de entrenamiento y de servicio. La clase optimiza principalmente $\mathcal{D}$, pero las restricciones siguen viniendo del modelo y de las GPU.
\end{knowledgebox}

\teachervoice{00:05:05--00:06:20, la ponente dice que las architecture/techniques se parecen cada vez más, por lo que, «con un presupuesto dado», son los datos los que marcan la diferencia entre modelos. La condición es given budget, no que a cualquier escala solo importen los datos.}

\subsection{Resumen de la sección}

Los open weights amplían el espacio de despliegue y personalización, pero la transparencia del entrenamiento sigue siendo la principal carencia. Por eso la clase reduce el preentrenamiento a una labor de ingeniería sujeta a la vez a restricciones de compute, serving, data, governance y evaluation, y StarCoder ofrece un caso de extremo a extremo que se puede verificar.
